\chapter{Threats to validity}
\label{ch:threats}

A result is only as strong as the ways it could be wrong, so they are listed here
rather than distributed through the text where they would be easier to skim past.

\section{Statistical}

\textbf{One seed per fixed-data cell.} The fixed-data sweep has only two surviving
cells, both sampled, and no training-seed replication. This is why the chapter
does not order tiny against small and does not claim a U-curve. The Banking77
transfer result has three seeds and a consistent direction; the AG News result has
three seeds but a paired interval that crosses zero.

\textbf{The exponent has no seed interval.} The matched-$D/N$ ladder is now
measured at four rungs, but at one training seed each. The per-point standard
errors ($\approx \IsoStderrMax$ bits/byte) describe held-out sampling, not
training variance, so $\alpha = \AlphaFit$ is a fitted line rather than a line
with an interval. That is enough to establish that the proxy improves
predictably with capacity; it is not enough to distinguish two rungs a few
percent apart.

\textbf{One seed per suite cell.} The nine-task, three-rung suite is
\SuiteTotalCells{} single-seed runs. Its binomial intervals describe
test-sampling error, not run-to-run variance, and the at-or-not-clearly-below
counts are descriptive. \SuiteReplicationCells{} cells were executed twice on the same
checkpoint and test draw to bound that variance directly: the repeats differ by at most
\SuiteReplicationMax{} points and flip up to \SuiteReplicationFlips{} of
2{,}000 predictions. Rung differences below that floor are not interpreted as
capability differences, and no equivalence claim rests on the counts.


\textbf{Fixed learning rate across a $47\times$ size range.} Comparability was
bought at the price of distance from each rung's own optimum. Part of every
size-related difference is optimisation, not capacity. The earlier 500-step,
\(5\times10^{-4}\) suite made that failure concrete: the mid-rung cells collapsed
for optimisation reasons, and the pre-registered sweep selected \(10^{-4}\) at
3{,}000 steps. The final suite uses that one setting for every task.



\section{Construct}

\textbf{Bits per byte is not accuracy.} It is a cheap, low-variance proxy. The
transfer matrix is the bridge to the quantity the thesis actually claims, and it
exists at one rung only. The pre-loss fixed-data curve is not available, so the
proxy is not used to establish a regime boundary in this build.

\textbf{AG News is a weak test and the result is inconclusive.} Banking77 has 77
fine-grained intents and shows a \TransferBankGain{} percentage-point gain over
random initialisation. AG News has four topics and a $25\%$ majority baseline; its
pooled gain is \TransferAgGain{} percentage points with a 95\% interval of
\TransferAgLo{} to \TransferAgHi{}, which is \TransferAgVerdict. A point estimate
inside a zero-crossing interval is not evidence of a gain, and it is not presented
as one.

\textbf{The suite is cross-harness.} The cited points come from different
tokenizers, pretraining corpora, model sizes, and evaluation harnesses. MASSIVE
intent's cited \SuiteCitedMassiveIntent\% is Laya's card, not a Jev number. A
point above a cited number is not a win, and interval overlap is not parity.

\textbf{A clean loss law and flat accuracy are one result, not a paradox.} The
matched-$D/N$ ladder lowers held-out bits/byte from \IsoTinyExh{} at 2.41M
parameters to \IsoLargeBpb{} at 113.47M, while the suite moves by
\SuiteLargeWinMin{} to \SuiteLargeWinMax{} points on \SuiteLargeWins{} cells,
loses \SuiteLargeLossAbs{} points on XNLI and leaves \SuiteLargeUnmoved{} cells
unmoved over
the same parameter range. The two measurements together rule out the reading
that decision accuracy at this data budget is waiting on more parameters. They
do not identify the true bottleneck --- a different packer, more task-relevant
pretraining data, or a different fine-tuning schedule could each move a cell ---
but they exclude capacity alone as sufficient for the four cells the commercial
claim needs.


\textbf{DAIR Emotion was a scoring artifact.} Its proportional test draw gives a
\SuiteEmotionMajority\% majority-class baseline, not a 1/6 chance rate. The
retired candidate-only values, \SuiteSmallPriorEmotion\% and
\SuiteMidPriorEmotion\%, were below that baseline; the dual-path values,
\SuiteSmallEmotion\% and \SuiteMidEmotion\%, are above it. The jump is a harness
correction, not a capability gain, and the earlier adverse reading is withdrawn.
The 33.7M cell is still \SuiteEmotionRungDrop{} points below the 14.4M cell, so
the correction does not create a scale result.

\textbf{Contamination.} The held-out region is the tail of a stream built from
public web text, and the transfer tasks are public datasets. The exact 40-byte
overlap audit found 3/500 AG News examples and 0/1,000 Banking77 examples; the
three AG matches scored 2/3, against 51.1\% on the unmatched 497. That is too
small and too easy to explain either transfer result, but it is a screen on exact
overlap, not a semantic deduplication.

\section{External}

\textbf{One corpus, one domain mix.} Byte-level bits per byte is corpus-dependent:
the bigram baseline is $3.767$ on the mixed stream and $4.117$ on a math-only
stream, so a baseline carried across a data change makes the gate meaningless. No
cross-corpus claim is made without refitting.

\textbf{No code in the corpus.} The planned source-code share is absent because its
dataset resolution hangs. For a thesis whose architecture argument rests partly on
code and audio sharing a byte space, this is a real omission, recorded in the
manifest rather than absorbed into the web-text share.

\textbf{Hardware-specific throughput.} The tok/s and VRAM figures are for one RTX
PRO 6000 Blackwell. They are measurements, not specifications, and they should not
be extrapolated to other accelerators without re-measuring.

\section{Methodological}

\textbf{A reclaimed artifact is not a result.} The first sandbox was reclaimed on
2026-10-28. The original AG News checkpoint \texttt{3cf421\ldots bed8}, the
exhaustive protocol replicates, the mid/large fixed-data cells and the context
runs were lost. A result that existed only in that VM is not
in this thesis. The tables mark those cells \emph{pending} and the prose gives the
reason; no old number is carried forward to make the table look complete.

\textbf{The checkpoint identity is amended.} The AG News cells use
\texttt{817e7f04\ldots0410}, retrained under the same protocol after the loss. It is
not bitwise-identical to the lost checkpoint. The Banking77 84.3\% row uses the
surviving pre-loss comparator \texttt{499fd68c\ldots ebdc}. The two rows are not a
single-checkpoint continuation, and the transfer table states that exception
rather than averaging across it.

\textbf{The queue is not a result.} A job that is running is not evidence. Sections
whose tables are pending are marked pending in prose, and the confidence summary in
\S\ref{sec:fixed-d} distinguishes what is established from what is scheduled.
